\section*{Overview}

Suffix automaton is the string structure I use when substring states matter more than lexicographic order. It stores all
substrings of one string in a DAG of only \(O(n)\) states and is especially strong when the string is processed online.

\section*{Problem-Driven Motivation}

Typical contest questions are:

\begin{itemize}
  \item how many distinct substrings does the string have,
  \item how many times does a pattern occur as a substring,
  \item what is the longest common substring with another string,
  \item can the text be extended online while maintaining substring information.
\end{itemize}

Trying to treat every substring separately is hopeless because there are \(O(n^2)\) of them. The useful compression is
to group substrings that behave the same with respect to ending positions.

\section*{Recognition Pattern}

Suffix automaton is a strong fit when:

\begin{itemize}
  \item the base string is scanned left to right,
  \item the queries are about substrings as automaton states, not about lexicographic suffix order,
  \item distinct-substring counting or occurrence aggregation is central,
  \item online extension is valuable.
\end{itemize}

If the problem is mainly about lexicographic order or suffix sorting, suffix array is often the cleaner tool.

\section*{Derivation}

The core abstraction is the \(\texttt{endpos}\) equivalence class:

\begin{quote}
Two substrings belong to the same state if they end at exactly the same set of positions in the full string.
\end{quote}

Each state stores:

\begin{itemize}
  \item \texttt{len[v]}: the maximum length in the class,
  \item \texttt{link[v]}: the suffix link to the largest proper suffix class,
  \item transitions by characters.
\end{itemize}

When appending one new character:

\begin{itemize}
  \item one new terminal state is created,
  \item suffix links are followed backward to add missing transitions,
  \item sometimes an existing state must be cloned because one old equivalence class actually splits into two classes.
\end{itemize}

The clone is the subtle part. It appears exactly when an old state would otherwise claim a shortest representative that is
too long, breaking the automaton invariant.

\section*{Worked Problem}

\paragraph{Problem.}
Build one text \(s\), then answer many pattern queries asking how many times each pattern appears in \(s\).

\paragraph{Why naive fails.}
Running a separate linear scan for every pattern is too slow when both the text and the query count are large.

\paragraph{Step 1: build the suffix automaton of the text.}
After processing the whole text, every substring corresponds to some path from the start state.

\paragraph{Step 2: propagate occurrence counts.}
While extending, each new terminal state contributes one end position. After building the automaton, propagate those
counts backward in decreasing \(\texttt{len}\) order along suffix links.

\paragraph{Why this propagation is correct.}
If state \(v\) has suffix link \(p\), then every occurrence represented by \(v\) also contributes to the suffix-class
represented by \(p\). So counts must flow from longer states to shorter suffix-link parents.

\paragraph{Step 3: answer pattern queries.}
Walk the pattern along automaton transitions. If a transition is missing, the pattern is absent. Otherwise the final
state's propagated count is exactly the number of occurrences.

\section*{Implementation Reasoning}

The meanings of the fields matter more than the code shape:

\begin{itemize}
  \item \texttt{len[v]} tells the largest length in the state,
  \item \texttt{len[link[v]] + 1 \ldots len[v]} is the full range of lengths represented by the state,
  \item \texttt{occ[v]} counts end positions only after the propagation pass.
\end{itemize}

The clone step copies transitions and the suffix link, but resets \texttt{occ} to zero because the clone is a structural
split, not a new end position.

\section*{Correctness Intuition}

The automaton stays minimal because states are exactly the end-position equivalence classes. Appending one character only
creates one genuinely new family of substrings: the suffixes ending at the new last position. Everything else is reused,
except when a previous state must be cloned to preserve correct suffix-link structure.

The formula
\[
\sum_{v \ne root} \bigl(len[v] - len[link[v]]\bigr)
\]
counts distinct substrings because each state contributes exactly the lengths that appear first in that state and in no
shorter suffix-link ancestor.

\section*{Complexity Analysis}

\begin{itemize}
  \item build: \(O(n)\) states and transitions for a fixed alphabet representation,
  \item occurrence propagation: \(O(n)\),
  \item one pattern existence or occurrence query: \(O(|pattern|)\).
\end{itemize}

\section*{Common Pitfalls}

\begin{itemize}
  \item Initializing clone states incorrectly.
  \item Forgetting that \texttt{occ} is not ready until the propagation pass is done.
  \item Using suffix automaton for lexicographic tasks that are simpler with suffix array.
  \item Hard-coding an alphabet that does not match the problem.
\end{itemize}

\section*{Variants and Failure Modes}

\begin{itemize}
  \item Longest common substring can be solved by walking a second string over the automaton.
  \item The suffix-link tree supports additional DP after the automaton is built.
  \item If the problem is about palindromes, palindromic tree is the right analog instead.
\end{itemize}

\section*{Practice Problems}

\begin{itemize}
  \item Count distinct substrings.
  \item Count occurrences of many patterns in one text.
  \item Longest common substring or repeated-substring tasks.
\end{itemize}

\section*{References}

\begin{itemize}
  \item \href{https://cp-algorithms.com/string/suffix-automaton.html}{cp-algorithms: Suffix Automaton}.
  \item \href{https://usaco.guide/adv/string-suffix}{USACO Guide: String Suffix Structures}.
  \item \href{https://codeforces.com/catalog}{Codeforces Catalog}.
\end{itemize}
